\documentclass[11pt,a4paper]{article}
\usepackage[margin=23mm,headheight=14pt,headsep=8.5mm,footskip=10mm]{geometry}
\usepackage{amsmath,amssymb,amsthm}
\usepackage{assets/tex/cm-bea}
\usepackage{xeCJK}
\setmainfont{Latin Modern Roman}
\setCJKmainfont{FandolSong-Regular.otf}[Script=CustomDefault,Path=assets/fonts/fandol/,BoldFont=FandolSong-Bold.otf,ItalicFont=FandolSong-Regular.otf,BoldItalicFont=FandolSong-Bold.otf]
\setCJKsansfont{FandolSong-Regular.otf}[Script=CustomDefault,Path=assets/fonts/fandol/,BoldFont=FandolSong-Bold.otf]
\xeCJKsetup{PunctStyle=quanjiao,CheckSingle=true}
\usepackage{setspace}
\setstretch{1.12}
\setlength{\parindent}{2em}
\setlength{\parskip}{3pt}
\usepackage{titlesec}
\titleformat{\section}{\centering\normalfont\Large\bfseries}{\thesection.}{0.6em}{}
\titleformat{\subsection}{\normalfont\large\bfseries}{}{0pt}{}
\titlespacing*{\section}{0pt}{1.8\baselineskip}{1.0\baselineskip}
\titlespacing*{\subsection}{0pt}{1.1\baselineskip}{0.5\baselineskip}
\usepackage{enumitem}
\setlist[enumerate]{leftmargin=2em,labelsep=0.5em,itemsep=3pt,parsep=0pt,topsep=6pt}
\providecommand{\tightlist}{\setlength{\itemsep}{3pt}\setlength{\parskip}{0pt}}
\usepackage{needspace}
\newenvironment{statementblock}{\par\Needspace{9\baselineskip}\noindent}{\par\nopagebreak[4]}
\newenvironment{proofblock}{\par\Needspace{4\baselineskip}\noindent}{\par}
\usepackage{fancyhdr}
\pagestyle{fancy}
\fancyhf{}
\fancyhead[L]{\small 数学手记}
\fancyhead[R]{\small 复变函数 · 第八章}
\fancyfoot[C]{\thepage}
\renewcommand{\headrulewidth}{0.4pt}
\usepackage[hidelinks,unicode,bookmarksnumbered=true,linktoc=all]{hyperref}
\usepackage{bookmark}
\bookmarksetup{depth=2}
\hypersetup{pdftitle={复变函数讲义：第八章},pdfauthor={Vincent Arqain},pdfsubject={调和函数与次调和函数}}
\renewcommand{\contentsname}{目录}
\setcounter{tocdepth}{1}
\setcounter{secnumdepth}{1}
\widowpenalty=10000
\clubpenalty=10000
\displaywidowpenalty=10000
\predisplaypenalty=10000
\interdisplaylinepenalty=10000
\raggedbottom

\usepackage{tabularx,booktabs,array}
\renewcommand{\thesection}{8.\arabic{section}}
\emergencystretch=2em

\usepackage{graphicx}
\usepackage{longtable}
\usepackage{calc}
\usepackage{float}
\floatplacement{figure}{H}
\renewcommand{\figurename}{图}
\makeatletter
\def\maxwidth{\ifdim\Gin@nat@width>\linewidth\linewidth\else\Gin@nat@width\fi}
\def\maxheight{\ifdim\Gin@nat@height>0.72\textheight0.72\textheight\else\Gin@nat@height\fi}
\makeatother
\setkeys{Gin}{width=\maxwidth,height=\maxheight,keepaspectratio}
\providecommand{\pandocbounded}[1]{#1}
\setcounter{tocdepth}{2}
\begin{document}
\begin{center}
{\LARGE\bfseries 复变函数讲义\par}
\vspace{8pt}{\large 第八章\quad 调和函数与次调和函数\par}
\vspace{6pt}{\small 依史济怀、刘太顺《复变函数》\quad 教材第 235—252 页\par}
\end{center}\vspace{8pt}
本章含4节讲义及全部 43 道习题原题。定义、定理和性质沿用教材的记号与编号；证明补足必要条件、退化情形与关键推导。覆盖 19 个教材编号条目，完整证明平均值性质、圆盘和半平面边值问题及次调和函数的基本性质。补齐一般上半连续情形的圆周平均证明，明确复积分收敛与负无穷值的处理；原题错误及必要条件独立校注。习题不附常规解答，原题中的疑义另列校勘。
\begingroup\setlength{\parskip}{0pt}\tableofcontents\endgroup\clearpage
\section{平均值公式与极值原理}
对应教材第 235---240 页。全章记
\(\mathbb D=B(0,1)\)、\(\mathbb H=\{z:\operatorname{Im}z>0\}\)，\(dA=dx\,dy\)。闭圆盘写作
\(\overline{B(a,r)}\)，\(G\Subset D\) 表示 \(\overline G\) 是 \(D\)
的紧子集。教材编号保持原样；标有``补充''的结论用于补全证明或说明适用条件。

\subsection*{定理索引}\label{ch08-s1-heading-1}

\hyperref[ch08-s1-heading-2]{调和性与局部全纯表示} · \hyperref[ch08-s1-heading-3]{定理
8.1.1} · \hyperref[ch08-s1-heading-4]{定理 8.1.2} · \hyperref[ch08-s1-heading-5]{定义
8.1.3} · \hyperref[ch08-s1-heading-6]{命题 8.1.4} · \hyperref[ch08-s1-heading-7]{定理
8.1.5} · \hyperref[ch08-s1-heading-8]{命题 8.1.6} · \hyperref[ch08-s1-heading-9]{补充}

\subsection*{调和性与局部全纯表示}\label{ch08-s1-heading-2}

实值函数 \(u\in C^2(D)\) 满足
\(\Delta u=u_{xx}+u_{yy}=0\)，称为调和函数。在单连通域中，闭微分形式
\(-u_y\,dx+u_x\,dy\) 有原函数 \(v\)，因为其闭合条件正是
\(-u_{yy}=u_{xx}\)。于是 \(v_x=-u_y,v_y=u_x\)，\(f=u+iv\)
全纯。一般域上可在每个小圆盘中作这个构造，故调和函数局部是全纯函数的实部，特别地属于
\(C^\infty(D)\)，且实解析。

共轭函数未必能在全域单值选取。例如 \(\log|z|\)
在环域上调和，其共轭是局部辐角，绕原点一周改变 \(2\pi\)。

若 \(\phi\) 全纯且 \(u\) 在其像域调和，链式法则与 Cauchy--Riemann
方程给出

\[
\Delta(u\circ\phi)=|\phi'|^2(\Delta u)\circ\phi=0.\tag{8.1.A}
\]

具体地，写 \(\phi=p+iq\)，二阶链式法则中含 \(\Delta p,\Delta q\)
的项为零，混合项的系数 \(\nabla p\cdot\nabla q=0\)，而
\(|\nabla p|^2=|\nabla q|^2=|\phi'|^2\)。此性质不要求 \(\phi\) 单射。

\Needspace{8\baselineskip}
\subsection{定理 8.1.1：平均值公式}\label{ch08-s1-heading-3}

若 \(u\) 在 \(B(a,R)\) 调和，则每个 \(0<r<R\) 满足

\[
u(a)=\frac1{2\pi}\int_0^{2\pi}u(a+re^{i\theta})\,d\theta.\tag{8.1.2}
\]

证明：在圆盘中选全纯函数 \(f\)，使 \(u=\operatorname{Re}f\)。Cauchy
积分公式给出

\[
\begin{aligned}
f(a)&=\frac1{2\pi i}\int_{|\zeta-a|=r}\frac{f(\zeta)}{\zeta-a}\,d\zeta\\
&=\frac1{2\pi}\int_0^{2\pi}f(a+re^{i\theta})\,d\theta.
\end{aligned}\tag{8.1.1}
\]

取实部即得 (8.1.2)。再乘以 \(2r\) 对 \(r\) 从 \(0\) 到 \(\rho\)
积分，利用极坐标换元，得到面积平均公式

\[
u(a)=\frac1{\pi\rho^2}\iint_{B(a,\rho)}u(z)\,dA(z).\tag{8.1.B}
\]

若 \(u\) 还连续到 \(\overline{B(a,R)}\)，由闭圆盘上的一致连续性，可令
\(r\uparrow R\)，故平均值公式也适用于外圆周。教材证明中全纯函数所在的空间应读为
\(H(B(a,R))\)。

\Needspace{8\baselineskip}
\subsection{定理 8.1.2：极值原理}\label{ch08-s1-heading-4}

域 \(D\) 上非常数的调和函数不能在内点取得全域最大值或最小值。

证明：若 \(u(a)=M=\sup_Du\)，则对每个以 \(a\) 为中心、闭包包含于 \(D\)
的小圆周，\(M-u\ge0\)
的平均值为零。连续非负函数积分为零必处处为零，故这些圆周上 \(u=M\)，从而
\(a\) 的一个邻域内 \(u=M\)。集合 \(E=\{z\in D:u(z)=M\}\)
因而既开又闭，且非空。由域的连通性，\(E=D\)，与非常数矛盾。对 \(-u\)
应用同一论证，得到最小值情形。

补充：非常数调和函数也不能取得局部极大值或极小值。若在一个小圆盘取得局部极大值，上述证明使它在该圆盘为常数；函数
\(u_x-iu_y\) 在全域全纯，在这个圆盘为零，由恒等定理在 \(D\) 为零，故
\(u\) 全域常数。

有界域 \(G\) 上的调和函数若连续到
\(\overline G\)，则最大值和最小值均可在 \(\partial G\)
取得。由紧性，极值在闭域上存在；如某个极值只在内点取得，前述原理使函数为常数，此时边界也取得同值。

\Needspace{8\baselineskip}
\subsection{定义 8.1.3：平均值性质}\label{ch08-s1-heading-5}

连续实值函数 \(u\) 对每个 \(\overline{B(a,r)}\subset D\) 都满足
(8.1.2)，称为具有平均值性质。这里连续性是定义的一部分，且要求所有这样的圆盘；只验证固定中心的一族圆周不够。

\Needspace{8\baselineskip}
\subsection{命题 8.1.4：平均值性质蕴含极值原理}\label{ch08-s1-heading-6}

具有平均值性质的非常数函数不能在域内取得全域最大值或最小值；若域有界且函数连续到闭域，极值可在边界取得。

证明：定理 8.1.2
关于全域极值的开闭集证明只用了连续性与平均值性质，所以逐句适用。此处先不假定函数可微。教材陈述中须补上``非常数''，否则常函数就是反例。§8.2
将证明平均值性质也保证调和性。

\Needspace{8\baselineskip}
\subsection{定理 8.1.5：Poisson 积分公式}\label{ch08-s1-heading-7}

若 \(u\) 在 \(B(0,R)\) 调和并连续到闭圆盘，则

\[
u(a)=\frac1{2\pi}\int_0^{2\pi}\frac{R^2-|a|^2}{|Re^{i\varphi}-a|^2}u(Re^{i\varphi})\,d\varphi,
\qquad |a|<R.\tag{8.1.9}
\]

证明：圆盘自同构

\[
\phi_a(z)=\frac{R^2(z-a)}{R^2-\overline a z}\tag{8.1.4}
\]

把 \(a\) 映到 \(0\)，且把半径 \(R\) 的圆周保持不变。由
(8.1.A)，\(u_1=u\circ\phi_a^{-1}\) 调和，\(u_1(0)=u(a)\)。对 \(u_1\)
使用边界平均值公式。令
\(\phi_a(Re^{i\varphi})=Re^{i\tau}\)，圆周上取对数微分并约去公共因子，得到

\[
\frac{d\tau}{d\varphi}=\frac{R^2-|a|^2}{|Re^{i\varphi}-a|^2}>0.\tag{8.1.7}
\]

这里正号说明保持圆周定向；\(\tau\) 应取连续增加的辐角，其总增量为
\(2\pi\)。将这个换元代入 \(u_1(0)\) 的平均值公式即得 (8.1.9)。写
\(a=re^{i\theta}\)，便得到教材的极坐标形式

\[
u(re^{i\theta})=\frac1{2\pi}\int_0^{2\pi}\frac{R^2-r^2}{R^2-2Rr\cos(\varphi-\theta)+r^2}u(Re^{i\varphi})\,d\varphi.\tag{8.1.10}
\]

\Needspace{8\baselineskip}
\subsection{命题 8.1.6：Poisson 核的三项性质}\label{ch08-s1-heading-8}

定义

\[
P(z,Re^{i\varphi})=\frac1{2\pi}\frac{R^2-|z|^2}{|Re^{i\varphi}-z|^2}.\tag{8.1.11}
\]

它严格为正，对 \(\varphi\in[0,2\pi]\) 的积分为 \(1\)，并且对圆内变量
\(z\) 调和。

证明：正性直接来自分子与分母。把 \(u\equiv1\) 代入
(8.1.9)，得质量归一化。固定 \(|\zeta|=R\)，有

\[
P(z,\zeta)=\frac1{2\pi}\operatorname{Re}\frac{\zeta+z}{\zeta-z}.\tag{8.1.C}
\]

分式在 \(|z|<R\) 全纯，故其实部调和。公式可由乘以
\(\overline\zeta-\overline z\) 直接核验。若使用弧长测度
\(|d\zeta|=R\,d\varphi\)，对应的核应是 \(P/R\)；本章统一采用角度测度
\(d\varphi\)。

Poisson
公式把边值的等权平均改为依赖内点的正权平均。越接近某个边界点，该点邻近弧段的权重越集中；下一节将把这一观察化为严格的边界极限证明。

\subsection*{补充：Harnack 估计与非负调和函数的控制}\label{ch08-s1-heading-9}

若 \(u\ge0\) 在 \(B(0,R)\) 调和，\(|z|=r<R\)，则

\[
\frac{R-r}{R+r}u(0)\le u(z)\le\frac{R+r}{R-r}u(0).\tag{8.1.D}
\]

证明：先取 \(r<\rho<R\)，在半径 \(\rho\) 的圆盘使用 Poisson
公式；核中的比值夹在 \((\rho-r)/(\rho+r)\) 与 \((\rho+r)/(\rho-r)\)
之间。边界函数非负，因此可逐点比较后积分，而其平均值为 \(u(0)\)。最后令
\(\rho\uparrow R\)。这也说明无须假设 \(u\) 连续到原圆盘边界。

沿域内连接两点的曲线取有限串重叠小圆盘，逐次应用
(8.1.D)，可把一处的非负函数值控制传递到另一处。对紧集再取有限覆盖，得到习题
8.1.9 所用的紧集一致比较常数。

\section{圆盘上的 Dirichlet 问题}
对应教材第 240---244 页。为区分边值与未知函数，本节以 \(g\)
记圆周上的已知实函数，以 \(h\)
记圆盘中的解。逐段连续指只有有限个第一类间断点，其余处连续；因圆周紧，这样的
\(g\) 有界。Dirichlet 问题要求 \(h\) 有界且调和，并在 \(g\) 的每个连续点
\(\zeta\) 满足从圆盘内部趋近的极限 \(h(z)\to g(\zeta)\)。

一般单连通域的 Riemann
映射只保证内部双全纯，不能单凭这一点就断言任意边界数据被一一搬到圆周。Jordan
域可以结合第七章的闭域同胚来处理；更一般边界须另检查边界对应或正则性。

\subsection*{定理索引}\label{ch08-s2-heading-1}

\hyperref[ch08-s2-heading-2]{定理 8.2.1} · \hyperref[ch08-s2-heading-3]{命题 8.2.2} ·
\hyperref[ch08-s2-heading-4]{定理 8.2.3} · \hyperref[ch08-s2-heading-5]{定理 8.2.4} ·
\hyperref[ch08-s2-heading-6]{补充}

\Needspace{8\baselineskip}
\subsection{定理 8.2.1：Poisson 积分的存在性与边界极限}\label{ch08-s2-heading-2}

设 \(g\) 在 \(|\zeta|=R\) 逐段连续，定义

\[
P[g](z)=\int_0^{2\pi}P(z,Re^{i\theta})g(Re^{i\theta})\,d\theta.\tag{8.2.1}
\]

则 \(P[g]\) 是圆盘中的有界调和函数，而且在每个连续点 \(\zeta_0\) 有

\[
\lim_{z\to\zeta_0,\ |z|<R}P[g](z)=g(\zeta_0).\tag{8.2.2}
\]

证明：设 \(|g|\le M\)。核的正性及质量为 \(1\) 给出
\(|P[g]|\le M\)。任取紧集 \(K\Subset B(0,R)\)，它到圆周的距离为正，故
Poisson 核关于 \(x,y\) 的一、二阶导数在 \(K\times\partial B(0,R)\)
一致有界。可以逐次在积分号下求导，由 \(\Delta_zP=0\) 得
\(\Delta P[g]=0\)。

给定 \(\varepsilon>0\)，由 \(g\) 在 \(\zeta_0\) 连续，选圆周上一小弧
\(I\)，使 \(|g(\zeta)-g(\zeta_0)|<\varepsilon\) 对 \(\zeta\in I\)
成立。利用核的质量，写成

\[
\begin{aligned}
|P[g](z)-g(\zeta_0)|
&\le\int_I P(z,\zeta)|g(\zeta)-g(\zeta_0)|\,d\theta\\
&\quad+\int_{\partial B(0,R)\setminus I}P(z,\zeta)|g(\zeta)-g(\zeta_0)|\,d\theta.
\end{aligned}\tag{8.2.A}
\]

第一项不超过 \(\varepsilon\)。令
\(d=\operatorname{dist}(\zeta_0,\partial B(0,R)\setminus I)>0\)；当
\(|z-\zeta_0|<d/2\) 时，远弧上 \(|\zeta-z|\ge d/2\)，所以第二项不超过

\[
\frac{8M}{d^2}(R^2-|z|^2)\longrightarrow0.\tag{8.2.B}
\]

这证明任意路径趋近的极限，绝不只限径向趋近。教材证明中远弧分母最终只需有一个固定正下界，不宜把其极限下界原封不动地写成邻域内的严格下界。

若 \(g\)
在整圆周连续，同一估计结合一致连续性，使解连续到整个闭圆盘。在间断点处，改变该点本身的赋值不改变积分；定理也不要求解取得这个任意赋值。

\Needspace{8\baselineskip}
\subsection{命题 8.2.2：有限例外边界点的极值估计}\label{ch08-s2-heading-3}

设 \(E\) 是 \(g\)
的连续点集合，\(m=\inf_Eg\)、\(M=\sup_Eg\)。任何满足这些边界极限的有界调和函数
\(h\) 均满足 \(m\le h\le M\)。

证明：设例外点为 \(\zeta_1,\ldots,\zeta_N\)。对 \(\varepsilon>0\) 定义

\[
v_\varepsilon(z)=M+\varepsilon\sum_{j=1}^N\log\frac{2R}{|z-\zeta_j|}.\tag{8.2.C}
\]

每项在圆内调和且非负，在对应的 \(\zeta_j\) 处趋于 \(+\infty\)。设
\(|h|\le A\)，取小半径 \(\delta\)，使靠近任一例外点的切割圆弧上
\(v_\varepsilon>A\)。在删去这些小圆帽后的域上，原外边界处
\(h\le M\le v_\varepsilon\)，新切割边界处也有
\(h<v_\varepsilon\)。这些域可能不连通，可逐个连通分支使用极值原理，故其中
\(h\le v_\varepsilon\)。

固定圆内一点 \(z\)，可将切割半径取到足够小以包含 \(z\)，再令
\(\varepsilon\downarrow0\)，得到 \(h(z)\le M\)。对 \(-h\)
作同样论证得下界。若没有例外点，直接使用闭圆盘极值原理。教材中 \(M\)
的上确界下标应与 \(m\) 一样取在 \(E\)，不是圆盘内部。

同一证明还给出：若有界调和函数在除有限个点外的边界均满足
\(\limsup h\le M\)，则圆内 \(h\le M\)。可先给上界添
\(\eta>0\)，用边界紧性找到邻近边界的统一控制，再令
\(\eta\downarrow0\)；因此不要求预先定义边界值。

\Needspace{8\baselineskip}
\subsection{定理 8.2.3：有界解的唯一性}\label{ch08-s2-heading-4}

逐段连续边值的圆盘 Dirichlet 问题有且仅有一个有界解，即 \(P[g]\)。

证明：存在性由定理 8.2.1。若 \(h_1,h_2\) 都是有界解，差 \(h_1-h_2\)
有界、调和，并在除有限个点外的边界趋于零。命题 8.2.2
对差给出上下界均为零，故二者相等。

有界性不可删去。例如

\[
h(z)=\operatorname{Re}\frac{1+z}{1-z}=\frac{1-|z|^2}{|1-z|^2}\tag{8.2.D}
\]

在单位圆盘内调和，在圆周除 \(1\) 外处处趋于零，却并非零函数，且在 \(1\)
附近无界。这与只有径向极限为零的习题情形也须区分。

\Needspace{8\baselineskip}
\subsection{定理 8.2.4：平均值性质刻画调和函数}\label{ch08-s2-heading-5}

域 \(D\) 上具有平均值性质的连续实值函数一定调和。

证明：任取 \(\overline{B(a,r)}\subset D\)，以圆周上的 \(u\) 为边值作
Poisson 积分 \(h\)。由定理 8.2.1，\(h\) 调和且连续到闭圆盘，并在边界等于
\(u\)。函数 \(u-h\) 连续并具有平均值性质，其边值为零。由命题 8.1.4
的极值原理，它在圆盘内恒为零。因此 \(u=h\) 在每个这样的圆盘上调和，故在
\(D\) 上调和。

这里并未预先对 \(u\) 求导；\(C^2\) 正则性是通过与 Poisson
积分一致而得到的。结合定理 8.1.1，调和性与连续函数的平均值性质等价。

\subsection*{补充：从 Fourier 模式看 Poisson 积分}\label{ch08-s2-heading-6}

对 \(0\le\rho<1\)，几何级数给出

\[
\frac{1-\rho^2}{1-2\rho\cos t+\rho^2}
=1+2\sum_{n=1}^{\infty}\rho^n\cos nt.\tag{8.2.E}
\]

证明：将 \(\sum_{n\ge0}\rho^ne^{int}=(1-\rho e^{it})^{-1}\)
取实部，再乘二减一；级数对固定 \(\rho<1\)
一致绝对收敛，允许逐项积分。因此边值 \(\cos n\theta\)、\(\sin n\theta\)
的解分别是 \((r/R)^n\cos n\theta\)、\((r/R)^n\sin n\theta\)。例如边值
\(2+3\cos2\theta-\sin\theta\) 的解为

\[
h(re^{i\theta})=2+3(r/R)^2\cos2\theta-(r/R)\sin\theta.\tag{8.2.F}
\]

这展示了边界振荡向内部衰减的速度，也给出检查常数、幂次和圆盘半径是否写对的简单方法。

\section{上半平面的 Dirichlet 问题}
对应教材第 244---247 页。把实轴与无穷远点合成扩充实轴；教材要求边值
\(g\) 在实轴只有有限个第一类间断点，且 \(\lim_{t\to+\infty}g(t)\) 与
\(\lim_{t\to-\infty}g(t)\) 均存在且有限。这保证 \(g\)
有界，两端极限不必相等。

\subsection*{定理索引}\label{ch08-s3-heading-1}

\hyperref[ch08-s3-heading-2]{定理 8.3.1} ·
\hyperref[ch08-s3-heading-3]{复积分形式的收敛校注} ·
\hyperref[ch08-s3-heading-4]{Schwarz--Christoffel 公式的调和函数证明}

\Needspace{8\baselineskip}
\subsection{定理 8.3.1：上半平面的 Poisson 公式}\label{ch08-s3-heading-2}

上述边值的唯一有界调和解为

\[
h(x+iy)=\frac1\pi\int_{-\infty}^{\infty}\frac{y}{(t-x)^2+y^2}g(t)\,dt,
\qquad y>0.\tag{8.3.4}
\]

证明：固定 \(z_0=x_0+iy_0\in\mathbb H\)，变换

\[
w=\phi(z)=\frac{z-z_0}{z-\overline{z_0}}\tag{8.3.A}
\]

把 \(\mathbb H\) 双全纯地映到 \(\mathbb D\)，把 \(z_0\) 映到
\(0\)，把扩充实轴映到圆周，无穷远点对应
\(1\)。搬运后的边值在圆周逐段连续，故由 §8.2 存在唯一有界解
\(h_1\)；\(h=h_1\circ\phi\) 即所求解。

实轴上 \(e^{i\theta}=\phi(t)\)，连续取圆周辐角，计算

\[
\begin{aligned}
i\,d\theta
&=\left(\frac1{t-z_0}-\frac1{t-\overline{z_0}}\right)dt\\
&=\frac{2iy_0}{(t-x_0)^2+y_0^2}\,dt.
\end{aligned}\tag{8.3.3}
\]

将它代入 \(h_1(0)\) 的圆周平均公式得到 (8.3.4)。由于 \(g\)
有界，积分绝对收敛；实核的积分为 \(1\)，所以
\(|h|\le\|g\|_\infty\)。在每个有限连续点，边界极限由圆盘结论及映射连续性得到。若两端极限相等，在无穷远点也取得该极限；若不等，那里是一个允许的边界间断点。

唯一性也由圆盘结论得到：两个有界解的差经变换后，在圆周除有限间断点及
\(1\)
外边值为零，有限例外点极值原理迫使差为零。此处不必给未知解预先规定无穷远点极限。若删去有界性，\(h(z)=\operatorname{Im}z\)
就是非零、实轴边值全为零的调和函数。

\subsection*{复积分形式的收敛校注}\label{ch08-s3-heading-3}

教材将 (8.3.4) 改写为
\(\operatorname{Re}\frac1{\pi i}\int_{\mathbb R}g(t)/(t-z)\,dt\)。逐点恒等式

\[
\operatorname{Re}\frac1{\pi i(t-z)}=\frac1\pi\frac{y}{(t-x)^2+y^2}\tag{8.3.B}
\]

确实成立，但在本节假设下，未取实部的复积分不一定收敛，甚至对称主值也可能不存在。例如
\(g(t)=1\) 对 \(t>0\)、\(g(t)=0\) 对 \(t<0\)，其单边积分含 \(\log T\)
发散项。

始终有效的解释是先对有限区间积分并取实部，再让两端独立趋于无穷；其结果由绝对收敛的实核积分给定。若需要真正的全纯函数，则采用正规化公式

\[
F(z)=\frac1{\pi i}\int_{\mathbb R}
\left(\frac1{t-z}-\frac{t}{1+t^2}\right)g(t)\,dt,
\qquad h(z)=\operatorname{Re}F(z).\tag{8.3.C}
\]

证明：括号中两项的 \(1/t\) 主部相消，在 \(z\) 的任意紧集上一致为
\(O(t^{-2})\)，故积分绝对且内闭一致收敛，积分内求导也合法，\(F\)
全纯。减去的项为实数，乘 \(1/i\)
后为纯虚数，不改变被积函数实部，故实部正是 (8.3.4)。

例如边值为区间 \((a,b)\) 的示性函数时，直接积分得

\[
h(x+iy)=\frac1\pi\left(\arctan\frac{b-x}{y}-\arctan\frac{a-x}{y}\right).\tag{8.3.D}
\]

这是该区间对点 \(z\) 的调和测度。它在区间内的边界趋于 \(1\)，区间外趋于
\(0\)；端点不属于连续点条件。

\subsection*{Schwarz--Christoffel 公式的调和函数证明}\label{ch08-s3-heading-4}

设 \(f:\mathbb H\to G\)
将上半平面双全纯地映到一个有界简单多边形内部，预顶点均为有限实数
\(a_1<\cdots<a_n\)，内角为
\(\alpha_k\pi\)，\(0<\alpha_k<2\)，并将无穷远点对应到一条边的内部。则

\[
f(z)=C\int_{z_*}^{z}\prod_{k=1}^n(\zeta-a_k)^{\alpha_k-1}\,d\zeta+C_1.\tag{8.3.6}
\]

证明：\(f'\) 无零点，故在单连通的 \(\mathbb H\) 上存在 \(L=\log f'\)
的全纯分支。记 \(v=\operatorname{Im}L\)。在实轴两相邻预顶点之间，Schwarz
反射把映射延到边界，\(v\) 的边界值是对应有向边的方向角，记为
\(\theta_k\)。这些角必须取同一个连续提升下的实数值，不能各自独立取主值。

还须先证明 \(v\) 有界，才能使用 Dirichlet 唯一性。第七章角点局部表示给出

\[
f(z)-w_k=(z-a_k)^{\alpha_k}A_k(z),\qquad A_k(a_k)\ne0,\tag{8.3.E}
\]

其中系数在适当邻域全纯；因此
\(L=(\alpha_k-1)\log(z-a_k)+\log B_k(z)\)，\(B_k\) 局部无零点。因
\(0<\arg(z-a_k)<\pi\)，其虚部在角点邻域有界。在无穷远点改用
\(\xi=-1/z\)，该点映到直边内部，反射后的映射导数非零，于是
\(f'(z)=c z^{-2}(1+O(1/z))\)，同样使 \(\operatorname{Im}L\)
有界。删去这些邻域后的闭集紧，\(L\) 连续，故 \(v\)
全域有界。这里仅使用第七章的局部反射与角点结论，未使用待证的整体乘积公式。

正向绕多边形经过第 \(k\) 个顶点时，有向切线转角为
\((1-\alpha_k)\pi\)，故

\[
\theta_k-\theta_{k-1}=(1-\alpha_k)\pi.\tag{8.3.F}
\]

对 \(0<\arg(z-a)<\pi\)，函数 \(\arg(z-a)/\pi\) 是实轴区间
\((-\infty,a)\) 的调和测度：在左侧趋于 \(1\)，右侧趋于
\(0\)，且有界、调和。于是

\[
v_0(z)=\theta_n+\sum_{k=1}^n(\alpha_k-1)\arg(z-a_k)\tag{8.3.G}
\]

在 \((a_j,a_{j+1})\) 上的边值为
\(\theta_n+\pi\sum_{k>j}(\alpha_k-1)=\theta_j\)。\(v\) 与 \(v_0\)
均为同一逐段常值边界问题的有界解，故相等。于是全纯函数

\[
L(z)-\sum_{k=1}^n(\alpha_k-1)\log(z-a_k)-i\theta_n
\]

虚部恒零。由 Cauchy--Riemann 方程它是实常数，指数化得到
\(f'=C\prod_k(z-a_k)^{\alpha_k-1}\)，再积分即得 (8.3.6)。积分路径位于
\(\mathbb H\)，各幂统一用上述对数分支定义，因单连通而路径无关。

教材在角度跳跃公式附近有下标及符号混淆。用 (8.3.F)
和每个区间的边值验算，可同时确定正确符号、系数及加法常数。这一证明的核心是先识别有界调和函数的边值，再由唯一性恢复导数的辐角。

\section{次调和函数}
对应教材第 247---252
页。平均值等式变为不等式后，函数可以在内点有低谷，甚至取
\(-\infty\)。因此本节把连续实值情形与一般上半连续情形分别说明，避免直接把连续边值的
Dirichlet 定理用于不连续函数。

\subsection*{定理索引}\label{ch08-s4-heading-1}

\hyperref[ch08-s4-heading-2]{定义 8.4.1} · \hyperref[ch08-s4-heading-3]{定理 8.4.2} ·
\hyperref[ch08-s4-heading-4]{定理 8.4.3} · \hyperref[ch08-s4-heading-5]{命题 8.4.4} ·
\hyperref[ch08-s4-heading-6]{命题 8.4.5} · \hyperref[ch08-s4-heading-7]{命题 8.4.6} ·
\hyperref[ch08-s4-heading-8]{命题 8.4.7} · \hyperref[ch08-s4-heading-9]{定理 8.4.8} ·
\hyperref[ch08-s4-heading-10]{本章联系与条件检查}

\Needspace{8\baselineskip}
\subsection{定义 8.4.1：次平均值性质与次调和函数}\label{ch08-s4-heading-2}

设 \(u:D\to\mathbb R\cup\{-\infty\}\)，\(u\not\equiv-\infty\)。若 \(u\)
上半连续，即

\[
\limsup_{z\to a}u(z)\le u(a),\tag{8.4.A}
\]

且对每个 \(\overline{B(a,r)}\subset D\) 有

\[
u(a)\le\frac1{2\pi}\int_0^{2\pi}u(a+re^{i\theta})\,d\theta,\tag{8.4.1}
\]

则称 \(u\)
次调和。上半连续函数在紧集上有有限上界，并且可测，因此积分可按扩充
Lebesgue 积分理解，允许为
\(-\infty\)，但不会出现正负无穷相减。调和函数显然次调和。

两个次调和函数的最大值仍次调和：在中心选取达到最大值的那一项，用它的次平均值不等式，再以两者的最大值放大被积函数。最大值的上半连续性也直接成立。特别地
\(u_N=\max(u,-N)\) 次调和，且 \(u_N\downarrow u\)。

\Needspace{8\baselineskip}
\subsection{定理 8.4.2：与调和函数的比较原理}\label{ch08-s4-heading-3}

设 \(u\) 是 \(D\) 上连续实值函数，则 \(u\) 次调和，当且仅当对每个
\(G\Subset D\) 及每个 \(h\in C(\overline G)\)、在 \(G\) 调和的函数，都有

\[
u\le h\text{ 在 }\partial G\quad\Longrightarrow\quad u\le h\text{ 在 }G.\tag{8.4.B}
\]

证明：先证必要性，并同时给出一般上半连续 \(u\) 也成立的方向。函数
\(w=u-h\) 上半连续且满足次平均值不等式。若 \(w\) 在闭域上最大值 \(M>0\)
于内点 \(a\) 取得，则对所有足够小的半径，圆周上的
\(w\le M\)，而次平均值要求其平均不小于 \(M\)，故平均必须等于
\(M\)。如果某圆周上一点的值小于
\(M\)，由上半连续性，在一段正长度弧上都有统一的严格亏量，积分便小于
\(M\)，矛盾。

因此 \(w=M\) 的集合在 \(G\) 中既开又闭，\(w\) 在域 \(G\) 恒为
\(M\)。取内点序列趋于任一边界点，上半连续性给出边界值至少为
\(M\)，与边界 \(w\le0\)
矛盾。这里函数在紧集上取最大值用的是上半连续性，不需要可微。

反之，给定闭圆盘，以连续边值 \(u\) 作 Poisson 解 \(h\)。由比较假设
\(u\le h\)，故

\[
u(a)\le h(a)=\frac1{2\pi}\int_0^{2\pi}u(a+re^{i\theta})\,d\theta.
\]

连续性又保证上半连续性，故 \(u\) 次调和。

上述必要性证明也给出强最大值原理：非常数次调和函数不能在域内取得全域最大值。但局部极大值可能存在，例如
\(u(z)=\max(0,\operatorname{Re}z)\)
次调和，在左半平面的任一点附近恒为零，因而取得非严格的局部极大值。

\Needspace{8\baselineskip}
\subsection{定理 8.4.3：圆周平均关于半径单调}\label{ch08-s4-heading-4}

若 \(u\) 在 \(\mathbb D\) 次调和，令 \(m(0)=u(0)\)，并定义

\[
m(r)=\frac1{2\pi}\int_0^{2\pi}u(re^{i\theta})\,d\theta,\qquad 0<r<1.
\]

则 \(m(r)\) 非降。这里结论包括允许 \(u=-\infty\) 的一般情形。

证明：先设 \(u\) 连续，取 \(0<r_1<r_2<1\)。以半径 \(r_2\) 圆周上的 \(u\)
为边值作调和函数 \(h\)。由比较原理 \(u\le h\)，故

\[
\begin{aligned}
m(r_1)&\le\frac1{2\pi}\int_0^{2\pi}h(r_1e^{i\theta})\,d\theta\\
&=h(0)=\frac1{2\pi}\int_0^{2\pi}u(r_2e^{i\theta})\,d\theta=m(r_2).
\end{aligned}\tag{8.4.C}
\]

对一般 \(u\)，不能直接声称存在一个连续到闭盘、边值等于 \(u\)
的调和函数。设外圆周为 \(K\)，在 \(K\) 上构造连续递减上逼近

\[
g_n(\zeta)=\max\left\{-n,\sup_{\eta\in K}\bigl(u(\eta)-n|\zeta-\eta|\bigr)\right\}.\tag{8.4.D}
\]

每个 \(g_n\) 有限且为 Lipschitz 函数，\(g_n\ge u\)，并随 \(n\)
递减。它在每一点下降到 \(u\)：若 \(u(\zeta)\)
有限，先用上半连续性控制近处的 \(u(\eta)\)，远处则由紧集上界和
\(-n|\zeta-\eta|\) 压低；若
\(u(\zeta)=-\infty\)，同一论证可用任意有限负阈值代替 \(u(\zeta)\)。若
\(u\) 在整个 \(K\) 为 \(-\infty\)，公式直接给出 \(g_n=-n\)。

令 \(h_n=P[g_n]\)。由上半连续版本的比较原理，圆内 \(u\le h_n\)，故
\(m(r_1)\le h_n(0)=\frac1{2\pi}\int g_n\)。对非负函数 \(g_1-g_n\)
用单调收敛定理，让 \(n\to\infty\)，得 \(m(r_1)\le m(r_2)\)，包括积分为
\(-\infty\) 的情况。\(r_1=0\)
则直接由定义。这样补上了教材从连续边值证明跨到一般次调和函数时的缺口。

\Needspace{8\baselineskip}
\subsection{命题 8.4.4：递增凸函数的复合}\label{ch08-s4-heading-5}

设 \(\varphi:\mathbb R\to\mathbb R\) 凸且非降，\(u\) 次调和。约定

\[
\varphi(-\infty)=\lim_{t\to-\infty}\varphi(t)\in[-\infty,\infty),
\]

则 \(\varphi\circ u\) 次调和。

证明：先以 \(u_N=\max(u,-N)\) 代替 \(u\)。在任一闭圆盘上，\(u_N\)
有上下界；由次平均值、\(\varphi\) 非降及 Jensen 不等式，

\[
\varphi(u_N(a))\le\varphi\left(\frac1{2\pi}\int u_N\right)
\le\frac1{2\pi}\int\varphi(u_N).\tag{8.4.E}
\]

这里 Jensen 不等式可由凸函数在平均值处的一条支撑直线
\(\varphi(t)\ge\varphi(t_0)+c(t-t_0)\)
积分得到；左右斜率的有序性保证这样的 \(c\) 存在。

由于 \(\varphi(u_N)\)
递减，且在圆周有共同有限上界，可减去该上界后用单调收敛传极限，得到
\(\varphi(u)\) 的次平均值不等式。有限凸函数在实轴连续，连同 \(-\infty\)
处的上述约定，保证复合上半连续；\(u\) 至少有一个有限值，故复合不恒为
\(-\infty\)。

\Needspace{8\baselineskip}
\subsection{命题 8.4.5：全纯函数产生的次调和函数}\label{ch08-s4-heading-6}

若 \(f\in H(D)\) 且 \(f\not\equiv0\)，则 \(\log|f|\) 和 \(|f|^p\) 对每个
\(p>0\) 都次调和。在零点处规定 \(\log0=-\infty\)。

证明：\(\log|f|\) 在非零点连续，在零点趋于
\(-\infty\)，故上半连续。若圆心是零点，次平均值不等式自动成立。否则固定
\(\overline{B(a,r)}\subset D\)，将 \(f\)
在这个闭圆盘内的所有零点按重数提出：

\[
f(z)=q(z)\prod_{j=1}^m(z-b_j),
\]

其中 \(q\) 在闭圆盘某邻域全纯且无零点，\(\log|q|\) 调和。对任意
\(b\)，基本圆周平均恒等式为

\[
\frac1{2\pi}\int_0^{2\pi}\log|a+re^{i\theta}-b|\,d\theta
=\log\max\{r,|b-a|\}.\tag{8.4.F}
\]

当 \(|b-a|\ne r\)
时，提出较大模长的因子，对剩余的非零全纯因子应用平均值公式即可。等号情形化为
\(\log r\) 加 \(\log|e^{i\theta}-1|\) 的平均；后者为零，因为
\(|e^{i\theta}-1|=2|\sin(\theta/2)|\)，并且
\(\int_0^\pi\log\sin t\,dt=-\pi\log2\)。此积分恒等式由
\(\sin2t=2\sin t\cos t\)、对称换元和
\(\int_0^\pi\log\sin t=2\int_0^{\pi/2}\log\sin t\)
联立得到；端点的对数奇性可积。

逐项取平均得到 Jensen 公式

\[
\frac1{2\pi}\int_0^{2\pi}\log|f(a+re^{i\theta})|\,d\theta
=\log|f(a)|+\sum_{|b_j-a|<r}\log\frac r{|b_j-a|}
\ge\log|f(a)|.\tag{8.4.3}
\]

圆周上的零点只贡献零项，其对数奇性由 (8.4.F)
已处理。因此不必暗中假设圆周无零点。最后取递增凸函数
\(\varphi(t)=e^{pt}\)，由命题 8.4.4 得 \(|f|^p\) 次调和，且这个证明对
\(0<p<1\) 同样有效。\(f\equiv0\) 时 \(|f|^p\equiv0\) 仍次调和，但
\(\log|f|\equiv-\infty\) 不满足本章排除恒负无穷的约定。

\Needspace{8\baselineskip}
\subsection{命题 8.4.6：全纯函数的积分平均}\label{ch08-s4-heading-7}

若 \(f\in H(\mathbb D)\)，\(0<p<\infty\)，则

\[
M_p(r,f)=\left(\frac1{2\pi}\int_0^{2\pi}|f(re^{i\theta})|^p\,d\theta\right)^{1/p}
\]

关于 \(r\) 非降。

证明：若 \(f\equiv0\) 显然成立；否则 \(|f|^p\) 由命题 8.4.5 次调和，定理
8.4.3 说明其圆周平均非降，再用 \(t\mapsto t^{1/p}\) 在 \([0,\infty)\)
递增即得。教材留作练习的这一步在此补齐。

\Needspace{8\baselineskip}
\subsection{命题 8.4.7：非负 Laplace
算子的弱最大值原理}\label{ch08-s4-heading-8}

设 \(u\in C^2(D)\) 为实值且 \(\Delta u\ge0\)，则对 \(G\Subset D\)，\(u\)
在 \(\overline G\) 的最大值可在 \(\partial G\) 取得。

证明：若 \(\Delta u>0\)，内点极大值处二阶导数满足
\(u_{xx}\le0,u_{yy}\le0\)，与严格正性矛盾。一般情形令
\(u_\varepsilon=u+\varepsilon|z|^2\)，则

\[
\Delta u_\varepsilon=\Delta u+4\varepsilon>0.\tag{8.4.G}
\]

因此
\(u(z)+\varepsilon|z|^2\le\max_{\partial G}(u+\varepsilon|\zeta|^2)\)。\(\overline G\)
有界，让 \(\varepsilon\downarrow0\) 得
\(u(z)\le\max_{\partial G}u\)。教材写成 \(\Delta|z|^2=1\) 的系数应修正为
\(4\)，不影响其扰动思路。

\Needspace{8\baselineskip}
\subsection{定理 8.4.8：二阶可微时的 Laplace 判据}\label{ch08-s4-heading-9}

对实值 \(u\in C^2(D)\)，次调和性等价于 \(\Delta u\ge0\)。

证明：若 \(\Delta u\ge0\)，则对比较原理中的 \(h\)，有
\(\Delta(u-h)\ge0\)。由命题 8.4.7，边界 \(u-h\le0\)
蕴含内部同样成立，再用定理 8.4.2 得次调和性。

反之，固定 \(a\)，对 \(u(a+re^{i\theta})\) 作二阶 Taylor
展开，对角度取平均。一次项、交叉项平均为零，\(\cos^2\theta\) 与
\(\sin^2\theta\) 的平均为 \(1/2\)，且 \(C^2\) 条件使余项对角度一致为
\(o(r^2)\)，故

\[
\frac1{2\pi}\int_0^{2\pi}u(a+re^{i\theta})\,d\theta-u(a)
=\frac{r^2}{4}\Delta u(a)+o(r^2).\tag{8.4.H}
\]

左端非负，除以 \(r^2\) 并令 \(r\downarrow0\) 得
\(\Delta u(a)\ge0\)。这也解释了次平均值不等式与二阶导数判据之间的精确关系。

\subsection*{本章联系与条件检查}\label{ch08-s4-heading-10}

平均值性质导出极值原理；Poisson 核建立 Dirichlet
解，再反过来证明平均值性质刻画调和性。将等式改为不等式，调和函数便成为次调和函数的比较对象，最后在
\(C^2\) 情形落到 \(\Delta u\ge0\)。

实际使用时应先检查：函数是否连续到边界，还是只给某些趋近方式的极限；解是否有界；无穷远点怎样处理；积分是否真的收敛；函数是否允许取
\(-\infty\)。这些条件决定唯一性、换元和极限步骤是否合法。第九章会将局部全纯性与映射结构推广到多个复变量。

\clearpage\section*{第八章习题}\phantomsection\addcontentsline{toc}{section}{第八章习题}
本页依教材印刷第 238---252 页收录四组共 43 道原题，题数分别为
19、6、3、15；保留原题的小问、提示和需要校核的表述。常规习题不附解答。疑误不在题干中静默改写，集中放在文末校注；其中的反例用于说明原题是否成立。

\subsection*{习题索引}\label{ch08-practice-heading-1}

\hyperref[ch08-practice-heading-2]{习题 8.1} · \hyperref[ch08-practice-heading-3]{习题 8.2} ·
\hyperref[ch08-practice-heading-4]{习题 8.3} · \hyperref[ch08-practice-heading-5]{习题 8.4} ·
\hyperref[ch08-practice-heading-6]{原题校注}

\Needspace{8\baselineskip}
\subsection{习题 8.1}\label{ch08-practice-heading-2}

\par\Needspace{5\baselineskip}
\noindent\textbf{1.} 证明：若 \(u\) 是域 \(D\) 上的调和函数，则
\(\dfrac{\partial^{j+k}u}{\partial x^j\partial y^k}\) 也是 \(D\)
上的调和函数（\(j,k=0,1,2,\ldots\)）。

\par\Needspace{5\baselineskip}
\noindent\textbf{2.} 证明：若 \(u\) 是域 \(D\) 上的调和函数，则
\(\dfrac{\partial u}{\partial z}\) 是 \(D\) 上的全纯函数。

\par\Needspace{5\baselineskip}
\noindent\textbf{3.} 设 \(f(z)\) 是域 \(D\) 上非常数的全纯函数，并且在 \(D\)
上不取零值。证明：

(1) \(\log|f(z)|\) 在 \(D\) 上调和；

(2) \(|f(z)|^p\) 在 \(D\) 上不调和，\(p>0\)。

\par\Needspace{5\baselineskip}
\noindent\textbf{4.} 证明：若域 \(D\) 上的调和函数列 \(\{u_n(z)\}\) 在 \(D\)
上内闭一致收敛于 \(u(z)\)，则 \(u\) 也是 \(D\) 上的调和函数，并且

\[
\left\{\frac{\partial^{j+k}u_n}{\partial x^j\partial y^k}\right\}
\longrightarrow\frac{\partial^{j+k}u}{\partial x^j\partial y^k}
\]

在 \(D\) 上内闭一致收敛。

\par\Needspace{5\baselineskip}
\noindent\textbf{5.} 证明：若 \(u\) 是域 \(D\) 上非常数的调和函数，则 \(u(D)\)
是开区间。

\par\Needspace{5\baselineskip}
\noindent\textbf{6.} 利用上题的结论证明：若 \(u\) 是域 \(D\)
上非常数的调和函数，则 \(u\) 不能在 \(D\)
内部取得极大值和极小值，自然不能在 \(D\) 内部取得最大值和最小值。

\par\Needspace{5\baselineskip}
\noindent\textbf{7.} 设 \(u(z)\) 在 \(B(0,1)\) 上调和，并且对任意
\(z_0\in\partial B(0,1)\) 成立 \(\lim_{r\to1}u(rz_0)=0\)，是否可以断言
\(u\equiv0\)？

（提示：举例说明答案是否定的。）

\par\Needspace{5\baselineskip}
\noindent\textbf{8.} 设 \(u\) 是 \(B(0,1)\) 上的非负调和函数，\(u(0)=1\)，给出
\(u(1/2)\) 的最佳估计。

\par\Needspace{5\baselineskip}
\noindent\textbf{9.} 设 \(D\) 是域，\(E\subset D\) 是紧集，\(a\in E\)。证明：若
\(u\) 是 \(D\) 上的非负调和函数，则存在仅与 \(a,E,D\) 有关的常数
\(\varepsilon\in(0,1)\)，使得

\[
\varepsilon u(a)\le u(z)\le\frac1\varepsilon u(a).
\]

\par\Needspace{5\baselineskip}
\noindent\textbf{10.}（Harnack 定理）设 \(\{u_n\}\) 是域 \(D\)
上单调增加的调和函数列。证明：若存在 \(z_0\in D\)，使得 \(\{u_n(z_0)\}\)
有界，则 \(\{u_n\}\) 在 \(D\) 上内闭一致收敛。

\par\Needspace{5\baselineskip}
\noindent\textbf{11.} 设 \(D\) 是由有限条光滑简单闭曲线围成的域，\(n\) 是
\(\partial D\) 指向外部的单位法向量场，\(u,v\) 在 \(\overline D\)
上调和。证明：

(1)
\(\displaystyle\int_{\partial D}\frac{\partial v(z)}{\partial n}|dz|=0\)；

(2)
\(\displaystyle\int_{\partial D}u(z)\frac{\partial v(z)}{\partial n}|dz|=\int_{\partial D}v(z)\frac{\partial u(z)}{\partial n}|dz|\)；

(3)

\[
\iint_D\left[\left(\frac{\partial u}{\partial x}\right)^2+\left(\frac{\partial u}{\partial y}\right)^2\right]dx\,dy
=\int_{\partial D}u(z)\frac{\partial u(z)}{\partial n}|dz|.
\]

\par\Needspace{5\baselineskip}
\noindent\textbf{12.} 设 \(u_1,\ldots,u_n\) 是域 \(D\) 上的调和函数。证明：若
\(|u_1(z)|+\cdots+|u_n(z)|\) 能在 \(D\) 内部取得极大值，则
\(u_1,\ldots,u_n\) 全部都是常值函数。

\par\Needspace{5\baselineskip}
\noindent\textbf{13.}（调和函数的 Hadamard 三圆定理）设
\(0<r_1<r_2<\infty\)，\(D=\{z:r_1<|z|<r_2\}\)，\(u\) 在 \(D\) 调和，在
\(\overline D\) 连续，\(A(r)=\max_{|z|=r}u(z)\)。证明：\(A(r)\) 在
\([r_1,r_2]\) 上是 \(\log r\) 的凸函数，即

\[
A(r)\le\frac{\log r_2-\log r}{\log r_2-\log r_1}A(r_1)
+\frac{\log r-\log r_1}{\log r_2-\log r_1}A(r_2).
\]

\par\Needspace{5\baselineskip}
\noindent\textbf{14.} 设 \(\varphi\in C^2((0,1))\)。证明：\(\varphi(|z|)\) 在
\(B(0,1)\setminus\{0\}\) 上调和，当且仅当存在实数 \(a,b\)，使
\(\varphi(|z|)=a\log|z|+b\)。

\par\Needspace{5\baselineskip}
\noindent\textbf{15.} 设 \(D\) 是凸域，\(u\) 是 \(D\) 上的调和函数。证明：若
\(u(z)>0\) 对所有 \(z\in D\) 成立，则存在 \(D\) 上的双全纯映射 \(f\)，使
\(\operatorname{Re}f'(z)=u(z)\)。

\par\Needspace{5\baselineskip}
\noindent\textbf{16.} 设 \(\gamma_1,\ldots,\gamma_n\) 为 \(n\)
条可求长曲线，\(u\) 是域 \(D\) 上的连续函数。若 \(u\) 在开集
\(D\setminus\bigcup_{k=1}^n\gamma_k\) 上调和，问 \(u\) 是否在 \(D\)
上调和？

（提示：举例说明答案是否定的。）

\par\Needspace{5\baselineskip}
\noindent\textbf{17.} 设 \(a_1,\ldots,a_n\) 是实数。证明：若
\(p(\theta)=\sum_{k=1}^n\cos k\theta\) 在 \([0,\pi]\) 上单调减少，则

\[
\sum_{k=1}^na_k\sin k\theta\ge0,\qquad\theta\in[0,\pi].
\]

（提示：注意 \(g(r,\theta)=\sum_{k=1}^nka_kr^k\sin k\theta\)
是单位圆盘中的调和函数。）

\par\Needspace{5\baselineskip}
\noindent\textbf{18.} 设 \(D\) 是域，\(a\in D\)。证明：若 \(u\) 在 \(D\)
上连续，在 \(D\setminus\{a\}\) 上调和，则 \(u\) 在 \(D\) 上调和。

\par\Needspace{5\baselineskip}
\noindent\textbf{19.}（Jensen 公式）设
\(f\in H(\overline{B(0,r)})\)，\(a_1,\ldots,a_m\) 是 \(f\) 在 \(B(0,r)\)
中的零点，重零点按重数重复计算，但 \(f(0)\ne0\)。证明：

\[
\log|f(0)|=-\sum_{k=1}^m\log\frac r{|a_k|}
+\frac1{2\pi}\int_0^{2\pi}\log|f(re^{i\theta})|\,d\theta.
\]

（提示：考虑
\(\displaystyle g(z)=f(z)\prod_{k=1}^m\frac{r^2-\overline{a_k}z}{r(z-a_k)}\)。）

\Needspace{8\baselineskip}
\subsection{习题 8.2}\label{ch08-practice-heading-3}

\par\Needspace{5\baselineskip}
\noindent\textbf{1.} 设 \(D\ne\mathbb C\) 为单连通域。证明：对任意
\(\zeta\in D\)，若 \(f_\zeta:D\to B(0,1)\)
双全纯，\(f_\zeta(\zeta)=0\)，\(f'_\zeta(\zeta)>0\)，则

(1) \(\log|f_\zeta(z)|=\log|f_z(\zeta)|\)，\(z,\zeta\in D\)；

(2)
\(\displaystyle\lim_{z\to z_0}\log|f_\zeta(z)|=0\)，\(z_0\in\partial D,\ \zeta\in D\)；

(3) 对固定 \(\zeta\in D\)，\(\log|f_\zeta(z)|-\log|z-\zeta|\) 作为 \(z\)
的函数在 \(D\) 调和。

\par\Needspace{5\baselineskip}
\noindent\textbf{2.} 设 \(D\) 是域，若
\(g:\overline D\times\overline D\to\mathbb R\) 满足

(1)
\(g(z,\zeta)=g(\zeta,z)\)，\((z,\zeta)\in\overline D\times\overline D\)；

(2) \(g\) 在
\((\overline D\times\partial D)\cup(\partial D\times\overline D)\)
上恒为零；

(3) 对固定 \(\zeta\in\overline D\)，\(g(z,\zeta)+\log|z-\zeta|\) 作为
\(z\) 的函数在 \(D\) 调和，在 \(\overline D\) 连续，

则称 \(g\) 为域 \(D\) 的 Green 函数。证明：异于 \(\mathbb C\)
的单连通域必有 Green 函数，并且唯一。

\par\Needspace{5\baselineskip}
\noindent\textbf{3.} 证明：若 \(g\) 是 \(B(0,R)\) 的 Green 函数，则

\[
P(z,Re^{i\theta})=-\frac1{2\pi}\lim_{r\to1}\frac{\partial g(z,rRe^{i\theta})}{\partial r},\qquad |z|<R.
\]

由此也可得到圆盘的 Poisson 积分公式和 Dirichlet 问题的解。

\par\Needspace{5\baselineskip}
\noindent\textbf{4.} 证明：\(B(0,1)\setminus\{0\}\) 的 Dirichlet 问题不可解。

\par\Needspace{5\baselineskip}
\noindent\textbf{5.}（Weierstrass 一致逼近定理）设 \(f\) 是 \(\partial B(0,R)\)
上的复连续函数。证明：对于任意 \(n\in\mathbb N\)，存在 \(z\) 的 \(n\)
次多项式 \(P_n(z)\) 和 \(\overline z\) 的 \(n\) 次多项式
\(Q_n(\overline z)\)，使 \(\{P_n(z)+Q_n(\overline z)\}\)
在圆周上一致收敛于 \(f\)。

\par\Needspace{5\baselineskip}
\noindent\textbf{6.} 设 \(\gamma_1,\gamma_2\) 为 \(\partial B(0,R)\)
上两段互余的开圆弧。求圆盘中的调和函数 \(u\)，使 \(\zeta\in\gamma_1\) 时
\(u(\zeta)=0\)，\(\zeta\in\gamma_2\) 时 \(u(\zeta)=1\)。并由此证明存在
\(f\in H(B(0,R))\)，使上述两弧上分别有 \(|f(\zeta)|=1\) 和
\(|f(\zeta)|=e\)。

\Needspace{8\baselineskip}
\subsection{习题 8.3}\label{ch08-practice-heading-4}

\par\Needspace{5\baselineskip}
\noindent\textbf{1.} 求上半平面 \(D\) 的 Green 函数 \(g\)，并证明：若

\[
P(z,t)=\frac1{2\pi}\lim_{s\to0}\frac{\partial g(z,t+si)}{\partial s},\qquad z\in D,\ t\in\mathbb R,
\]

则 \(\displaystyle u(z)=\int_{-\infty}^{\infty}P(z,t)f(t)\,dt\)
是以上半平面实轴上的 \(f(t)\) 为边值的 Dirichlet 问题的解。

\par\Needspace{5\baselineskip}
\noindent\textbf{2.} 设 \(D\) 是由简单闭曲线围成的单连通域。证明：\(D\) 的
Dirichlet 问题可解，并且解唯一。

\par\Needspace{5\baselineskip}
\noindent\textbf{3.} 求解上半单位圆盘 Dirichlet 问题的具体公式。

\Needspace{8\baselineskip}
\subsection{习题 8.4}\label{ch08-practice-heading-5}

\par\Needspace{5\baselineskip}
\noindent\textbf{1.} 证明：若 \(u\) 在域 \(D\) 次调和，则对任意 \(z_0\in D\)，

\[
\limsup_{z\to z_0}u(z)=u(z_0).
\]

\par\Needspace{5\baselineskip}
\noindent\textbf{2.} 证明：域 \(D\) 上非常数的次调和函数不能在 \(D\)
内部取得最大值。举例说明它能在 \(D\) 内部取得极大值。

\par\Needspace{5\baselineskip}
\noindent\textbf{3.} 证明：若 \(u\) 是域 \(D\) 上非常数的连续次调和函数，则
\(u(D)\) 是右开的区间。

\par\Needspace{5\baselineskip}
\noindent\textbf{4.} 证明：若 \(u\) 是 \(B(0,R)\) 上的连续次调和函数，则

\[
m(r)=\frac1{\pi r^2}\iint_{B(0,r)}u(z)\,dx\,dy
\]

是 \((0,R)\) 上的增加函数。

\par\Needspace{5\baselineskip}
\noindent\textbf{5.} 举例说明，存在域 \(D\) 上的上半连续函数 \(u\)，它不能在
\(D\)
内部取得最大值，但不是次调和函数。这说明最大值原理不是次调和函数的特征性质。

\par\Needspace{5\baselineskip}
\noindent\textbf{6.} 证明：域 \(D\) 上有限个次调和函数的和仍为次调和函数。

\par\Needspace{5\baselineskip}
\noindent\textbf{7.} 证明：若 \(\{u_n\}\) 是域 \(D\) 上单调减少的次调和函数列，且
\(\lim_{n\to\infty}u_n=u\not\equiv-\infty\)，则 \(u\) 次调和。

\par\Needspace{5\baselineskip}
\noindent\textbf{8.} 设 \(u_n\) 是域 \(D\) 上的次调和函数列。证明：若
\(\sup_{n\ge1}u_n=u\not\equiv\infty\)，且 \(u\) 是 \(D\)
上的上半连续函数，则 \(u\) 次调和。

\par\Needspace{5\baselineskip}
\noindent\textbf{9.} 证明：若 \(u\) 在域 \(D\) 调和，则对每个
\(p\ge1\)，\(|u|^p\) 次调和。

\par\Needspace{5\baselineskip}
\noindent\textbf{10.}（次调和函数的 Hadamard 三圆定理）设
\(0<r_1<r_2<\infty\)，\(D=\{z:r_1<|z|<r_2\}\)，\(u\) 在 \(D\) 次调和，在
\(\overline D\) 连续，\(A(r)=\max_{|z|=r}u(z)\)。证明：\(A(r)\) 是
\(\log r\) 的凸函数，即

\[
A(r)\le\frac{\log r_2-\log r}{\log r_2-\log r_1}A(r_1)
+\frac{\log r-\log r_1}{\log r_2-\log r_1}A(r_2).
\]

\par\Needspace{5\baselineskip}
\noindent\textbf{11.} 设 \(D,G\) 是域，\(u\) 在 \(G\) 次调和。证明：若
\(\varphi:D\to G\) 全纯，则 \(u\circ\varphi\) 在 \(D\) 次调和。

\par\Needspace{5\baselineskip}
\noindent\textbf{12.} 设 \(D\) 是凸域，\(f\in H(D)\)。证明：若
\(\operatorname{Re}(\overline z f(z))\) 在 \(D\) 次调和，则 \(f\)
是常值函数或是 \(D\) 上的双全纯映射。

\par\Needspace{5\baselineskip}
\noindent\textbf{13.} 设 \(\gamma_1,\ldots,\gamma_n\) 为 \(n\)
条可求长曲线，\(u\) 是域 \(D\) 上的连续函数。若 \(u\) 在开集
\(D\setminus\bigcup_{k=1}^n\gamma_k\) 上次调和，问 \(u\) 是否在 \(D\)
上次调和？

（提示：举例说明答案是否定的。）

\par\Needspace{5\baselineskip}
\noindent\textbf{14.} 设 \(v\) 是有界域 \(D\) 上的次调和函数，\(u\) 是 \(D\)
上的调和函数。证明：若存在 \(z_1,\ldots,z_n\in\partial D\)，使

\[
\begin{aligned}
\limsup_{z\to z_0}[v(z)-u(z)]&\le0,
&&z_0\in\partial D\setminus\{z_1,\ldots,z_n\},\\
\limsup_{z\to z_k}[v(z)-u(z)]&<\infty,
&&k=1,\ldots,n,
\end{aligned}
\]

则 \(v(z)\le u(z)\) 对所有 \(z\in D\) 成立。

\par\Needspace{5\baselineskip}
\noindent\textbf{15.}（两常数定理）设 \(\gamma,\Gamma\)
是单位圆周上两段不相交的开圆弧，长度分别为 \(2\alpha\pi\) 和
\(2(1-\alpha)\pi\)，\(f\) 是单位圆盘上的有界全纯函数。证明：若 \(m,M>0\)
且

\[
\begin{aligned}
\limsup_{z\to z_0}|f(z)|&\le m,&&z_0\in\gamma,\\
\limsup_{z\to z_0}|f(z)|&\le M,&&z_0\in\Gamma,
\end{aligned}
\]

则 \(|f(0)|\le m^\alpha M^{1-\alpha}\)。由此可得到 \(|f(z)|\)
的什么估计？

（提示：参考习题 8.2 的第 6 题。）

\Needspace{8\baselineskip}
\subsection{原题校注与适用条件}\label{ch08-practice-heading-6}

\textbf{8.1.9 的量词。} 不等式应对 \(z\in E\)
成立；原题显示公式未再次写出这一范围。常数不能要求对整个任意域 \(D\)
同时适用。

\textbf{8.1.11 的边界正则性。} ``在 \(\overline D\)
上调和''通常表示在其某开邻域调和；也可采用在 \(D\)
调和且具有足够边界可微性的表述。不能仅凭在开域调和，就默认法向导数在边界存在。

\textbf{8.1.12 的结论不成立。} 在 \(D=B(0,1)\) 取
\(u_1(z)=2+\operatorname{Re}z\)、\(u_2(z)=2-\operatorname{Re}z\)，则两者都为正且非常数，但
\(|u_1|+|u_2|\equiv4\)，处处取得极大值。一个有效修订是把绝对值之和改为平方和：因为
\(\Delta\sum_j u_j^2=2\sum_j|\nabla u_j|^2\)，局部极大值先使次调和平方和在一邻域为常数，再迫使各梯度为零，最后用调和函数的唯一延拓得到全域常数。原题未标记为成立。

\textbf{8.1.17 漏了系数。} 原书写
\(p(\theta)=\sum\cos k\theta\)，其假设与 \(a_k\) 无关。取 \(n=1,a_1=-1\)
即反驳原结论。结合提示，应修订为
\(p(\theta)=\sum a_k\cos k\theta\)；原稿题干保留原式。

\textbf{8.1.19 的边界零点。} \(H(\overline{B(0,r)})\)
解释为在闭圆盘某邻域全纯。圆周零点导致可积的对数奇性，不能套用要求连续边值的平均值公式；§8.4
的 (8.4.F) 已补足处理。

\textbf{8.2.1 的对角线。} 第 (1) 小问在 \(z=\zeta\) 时两侧为
\(-\infty\)，按扩充值理解；第 (3) 小问在 \(z=\zeta\) 处应作可去延拓。

\textbf{8.2.2 的 Green 函数定义需修订。} 极点必须固定在
\(\zeta\in D\)，不能放到
\(\partial D\)；函数在内部对角线上有正对数奇性，也不能宣称处处取有限实数。通常在
\(z\ne\zeta\) 处定义，要求 \(g(z,\zeta)+\log|z-\zeta|\)
可去延拓为调和函数，并在有限边界趋于零。有界域中可用极值原理证明唯一性；对无界域还须指定无穷远条件或采用最小正
Green 函数的规范。

仅有有限边界零值不足以保证无界域的唯一性：在上半平面，标准核
\(g_0(z,\zeta)=\log\left|(z-\overline\zeta)/(z-\zeta)\right|\) 与
\(g_0(z,\zeta)+c\operatorname{Im}z\operatorname{Im}\zeta\)（\(c>0\)）都满足对称性、有限边界零值与同样的对数主部，却不同。故不把原题字面上的定义和无条件唯一性标记为正确。

\textbf{8.2.4 的``不可解''。}
指存在边界数据无解，不是每组数据都无解。圆周赋值零、孤立边界点 \(0\)
赋值 \(1\) 就不可能有连续至全边界的调和解；零边值当然仍有零解。

\textbf{8.2.5 的次数。} 一致逼近通常采用次数``不超过
\(n\)''的多项式；若要求恰为 \(n\)
次，可添加趋零的高次项，不改变一致收敛。并非对固定 \(n\)
已能任意精度逼近。

\textbf{8.2.6 与 8.3.1---3 的边值。}
取值指从域内趋近的极限；需沿用本章连续或逐段连续边值及有界解的约定。半平面
Green 核取上述标准规范，法向导数中的 \(s\to0\) 为 \(s\downarrow0\)。

\textbf{8.4.1、8.4.14、8.4.15 的横线。} 原书极限符号上方的横线表示
\(\limsup\)，不是普通极限，已按其数学含义排写。

\textbf{8.4.4 的``增加''。} 应理解为非降，不是严格递增，常函数就是反例。

\textbf{8.4.8 的有限值约定。} 按本章定义，须保证 \(u(z)<+\infty\)
对每点成立；``不恒为
\(+\infty\)''本身不能替代这一值域条件。上半连续性应按值域
\(\mathbb R\cup\{-\infty\}\)
理解，不能先允许局部为正无穷再直接调用定义。

\textbf{8.4.11 的恒负无穷例外。} 因本章排除恒为 \(-\infty\) 的函数，需补
\(u\circ\varphi\not\equiv-\infty\)。例如
\(u(w)=\log|w|\)、\(\varphi\equiv0\) 时复合恒为
\(-\infty\)。对非常数全纯映射，或常值落在 \(u\)
有限处的映射，结论按通常定义成立；若把恒负无穷也纳入广义次调和函数，原句可直接保留。


\end{document}
